\chapter{The resolution ladder and the release certificate}\label{ch:ladder}
\chapterintro{A specification is solved on a ladder of grids, each twice as fine per axis as the one below. This chapter
explains how the rungs are chosen, how one rung starts from the design of the one below, the three warm release
policies and what each can and cannot do, and the six predicates of the release certificate.}

\section{The rungs}\label{sec:rungs}
\index{resolution ladder}\index{rung}
The ladder is built from the specification's base mesh $h$. Candidate element sizes are $h\cdot 2^k$ for integer
$k$; a candidate is kept when its grid divides the domain into whole elements and has between \LadderMin\ and
\LadderMax\ elements. Consecutive rungs must differ by exactly a factor of two per axis, so that every rung can start
from the one below. The four finest rungs are named \emph{sketch}, \emph{preview}, \emph{refine} and \emph{release};
any coarser ones are named \code{coarse0}, \code{coarse1}, \ldots{} and are not shown in the workspace.

\begin{figure}[h]
\centering
\begin{tikzpicture}[font=\sffamily\footnotesize,
  rung/.style={draw=#1,fill=#1!10,line width=0.7pt,rounded corners=1.5pt,minimum height=9mm,align=center},
  arr/.style={-{Stealth[length=2mm]},line width=0.6pt,color=guigrey}]
  \node[rung=guigrey,minimum width=20mm] (s) {\textbf{Sketch}\\144k};
  \node[rung=guigrey,minimum width=24mm,right=7mm of s] (p) {\textbf{Preview}\\1.15M};
  \node[rung=guigrey,minimum width=28mm,right=7mm of p] (r) {\textbf{Refine}\\9.2M};
  \node[rung=okgreen,minimum width=34mm,right=7mm of r] (R) {\textbf{Release}\\73.7M, spill guard on};
  \draw[arr] (s) -- node[above,font=\scriptsize\sffamily]{warm} (p);
  \draw[arr] (p) -- node[above,font=\scriptsize\sffamily]{warm} (r);
  \draw[arr] (r) -- node[above,font=\scriptsize\sffamily]{warm} (R);
  \node[below=1.5mm of s,font=\scriptsize\sffamily,color=muted] {seconds, cold};
  \node[below=1.5mm of p,font=\scriptsize\sffamily,color=muted] {seconds};
  \node[below=1.5mm of r,font=\scriptsize\sffamily,color=muted] {seconds};
  \node[below=1.5mm of R,font=\scriptsize\sffamily,color=muted] {seconds to minutes};
\end{tikzpicture}
\caption{The ladder of the motor mount b06: element counts per rung. Each rung has eight times the elements of the one
below and starts from its design. Only the release rung runs with spill monitoring, so only it can be certified.}
\label{fig:ladder}
\end{figure}

The release rung is the finest grid on the ladder, not the specification's $h$. When a brief asks for elements of at
most 1~mm, a release at 0.25~mm satisfies it; \tool{ladder} says so in a note. Asking \tool{solve} for a rung and an
\code{h} that disagree is an error, so an agent cannot accidentally solve the refine grid under the release name.

\paragraph{Solver choice.} Grids up to \MgSwitch\ elements use the FP64 multigrid solver \code{mg64}; larger grids use
\code{mg\_ir}, a mixed-precision iterative-refinement variant that needs less memory. If \code{mg\_ir} fails on a grid
of at most \MgRetry\ elements, the worker retries once on \code{mg64}.

\paragraph{Typical times.} Table~\ref{tab:rungtimes} gives the latency study's numbers on the RTX~4090.

\begin{table}[h]
\caption{Backend timings over the 14 released tasks of the session study, each measured through the running backend
\claim{C6.4, C6.5}. Warm means started from the design of the rung below or, for the sketch, from the previous sketch
after an edit.}\label{tab:rungtimes}
\begin{fixtable}{@{}>{\raggedright\arraybackslash}p{0.38\linewidth}X@{}}
Interaction & Range over tasks \\ \midrule
Edit $\rightarrow$ validation, 95th percentile & 1.2--10.5 ms \\
Edit $\rightarrow$ warm sketch, 95th percentile & 0.59--2.40 s \\
Cold sketch of a new specification & 2.5--6.7 s (4 of 14 within 3 s); first density frame after about 0.1 s \\
Release, warm (tail policy), compute & 3.1--56.4 s, all 14 certified \\
\end{fixtable}
\end{table}

\section{Warm starts}\label{sec:warm}
\index{warm start}
When a solve has a suitable earlier design, it starts from it:
\begin{enumerate}
  \item The coarse density is upsampled by a factor of two in each direction (piecewise constant).
  \item One pass of the fine rung's density filter smooths it, so boundary elements take intermediate values. Without
  this step an upsampled 0/1 design is saturated under the sharp final projection, and no small move can change it.
  \item A constant density offset, found by bisection, makes the projected design volume equal the target. Keep-in and
  keep-out regions voxelize differently at the two resolutions, so the upsampled volume is not exactly the target.
\end{enumerate}

The warm source is chosen automatically. A solve on a rung uses the design of the same specification one doubling
coarser, if there is one (even a solve still running, behind which it queues). After an edit, when no coarser design of
the new specification exists, it uses the latest design on the same grid. Without either, the solve is cold.

\section{Release policies}\label{sec:policies}
\index{release policy}\index{tail schedule}\index{long re-opening}\index{short re-opening}
After the warm start, the optimizer runs a continuation schedule. Table~\ref{tab:schedules}, generated from the code,
lists the four schedules. Each row runs until the given iteration and sets the SIMP penalty $p$, the projection
sharpness $\beta$, the move limit and the filter radius $r_\mathrm{min}$ in element units. All runs stop at
convergence or at \MaxOuter\ iterations.

\begin{table}[htb]
\caption{Continuation schedules (from \code{solver.py} and \code{escalation.py}).}\label{tab:schedules}
% Generated by docs/guidebook/gen_tables.py from the code. Do not edit by hand.
\begin{fixtable}{@{}>{\raggedright\arraybackslash}p{0.21\linewidth}rrrrr>{\raggedright\arraybackslash}X@{}}
Policy & Until it. & $p$ & $\beta$ & Move & $r_\mathrm{min}$ & Starts from \\ \midrule
Cold (reference) & 15 & 1.5 & 1 & 0.20 & 1.50 & uniform density \\
 & 40 & 3.5 & 4 & 0.15 & 1.35 &  \\
 & 65 & 4.5 & 16 & 0.08 & 1.25 &  \\
 & 120 & 4.5 & 32 & 0.05 & 1.20 &  \\
\addlinespace[2pt]
Tail & end & 4.5 & 32 & 0.05 & 1.20 & warm from the rung below (default) \\
\addlinespace[2pt]
Short re-opening & 4 & 4.5 & 8 & 0.10 & 1.20 & warm, same grid after an edit \\
 & 10 & 4.5 & 16 & 0.08 & 1.20 &  \\
 & end & 4.5 & 32 & 0.05 & 1.20 &  \\
\addlinespace[2pt]
Long re-opening & 25 & 3.5 & 4 & 0.15 & 1.35 & warm from the rung below, \textsf{Optimize release} \\
 & 50 & 4.5 & 16 & 0.08 & 1.25 &  \\
 & end & 4.5 & 32 & 0.05 & 1.20 &  \\
\addlinespace[2pt]
\end{fixtable}

\end{table}

\begin{description}
  \item[Tail] (\gui{Release}, the default warm start from a coarser rung). Only the final phase: $p = 4.5$,
  $\beta = 32$, move 0.05. It takes about 12 iterations and mostly \emph{certifies} the upsampled design on the fine grid
  rather than improving it.
  \item[Long re-opening] (\gui{Optimize release}, or \code{optimize=True} in \tool{solve}). The cold schedule from its
  second phase on: $p = 3.5$, $\beta = 4$, move 0.15, then sharpening to the final phase. The projection starts far from
  saturation, so boundaries move and members finer than the coarse rung can form. About 60 iterations.
  \item[Short re-opening] (automatic after an edit, and for clarification previews). $\beta$ steps 8, 16, 32 with
  larger moves, so a design on the same grid can follow a changed specification.
  \item[Cold.] The full schedule from uniform density: what a solve does without a warm source, and the reference in
  the paper's comparisons.
\end{description}

\begin{warnbox}[A tail release is not a fine-grid optimization]
A sharply projected upsampled design can be saturated away from the projection threshold; continuing only the final
phase then changes essentially no solid or void states. A certified tail release is an acceptance of the transferred
coarse design on the fine grid, which is useful, but it is not an optimization from a uniform fine-grid start. In the
paper's comparison the binary design changed essentially not at all after a tail release. Use \gui{Optimize release} for the final design.
\end{warnbox}

\begin{evidencebox}[Release policies on five problems (paper, Table 8)]
At their release grids of 25.6--73.7 million elements, every final result below was certified. $\Delta C$ is the
compliance relative to a cold run on the same grid; compute is the share of the cold run's release-rung compute
(earlier rungs excluded) \claim{C3.1}.

\smallskip
{\centering\sffamily\footnotesize
\begin{tabular}{@{}lrrrrrr@{}}
\toprule
Problem & Elements & Cold & Tail $\Delta C$ & Tail compute & Long $\Delta C$ & Long compute \\ \midrule
Cantilever & 65.5M & 118 s & +0.28\% & 7\% & $-$1.20\% & 69\% \\
Bridge & 65.5M & 249 s & +1.39\% & 11\% & $-$1.19\% & 51\% \\
Torsion & 65.5M & 121 s & +0.86\% & 6\% & +0.07\% & 59\% \\
L-bracket & 25.6M & 251 s & +0.00\% & 3\% & $-$3.06\% & 51\% \\
Motor mount & 73.7M & 423 s & +1.22\% & 8\% & +0.03\% & 61\% \\
\bottomrule
\end{tabular}\par}

\smallskip
The long re-opening reached a different local optimum from the cold run (binary overlap 0.86--0.97); a lower compliance
is evidence of a different trajectory, not of a global optimum. The bridge's cold reference needed a different solver
configuration after two failures. A sixth problem, a duct bridge, gave no usable comparison and is excluded.
\end{evidencebox}

\section{The release certificate}\label{sec:certificate}
\index{certificate!predicates}
Every finished solve carries a certificate, created inside the solver from its own evidence. The tool contract can read
it but has no way to write it. The certificate is the conjunction of six predicates:
\begin{equation}
  A_\mathrm{num} = A_\mathrm{conv}\land A_\mathrm{res}\land A_\mathrm{no\,spill}\land A_\mathrm{monitor}\land
  A_\mathrm{vol}\land A_\mathrm{passive}.
\end{equation}
The final state solve is accepted on the true relative residual, recomputed in FP64:
\begin{equation}
  r_\mathrm{true} = \frac{\lVert f - K u\rVert_2}{\lVert f\rVert_2} \le \tau ,
\end{equation}
where $\tau$ is the solve's configured tolerance (normally $10^{-5}$ for optimization in the paper's experiments).
Table~\ref{tab:checks}, generated from \code{solver.py} and the result card's labels, lists the six predicates.

\begin{table}[htb]
\caption{The six certificate checks. A result is certified only when all six pass.}\label{tab:checks}
% Generated by docs/guidebook/gen_tables.py from the code. Do not edit by hand.
\begin{fixtable}{@{}>{\raggedright\arraybackslash}p{0.20\linewidth}>{\raggedright\arraybackslash}p{0.19\linewidth}l>{\raggedright\arraybackslash}X@{}}
Check (record key) & Result-card label & Predicate & Passes when \\ \midrule
\code{converged} & converged & $A_\mathrm{conv}$ & The optimizer met its stopping criterion within the iteration cap. \newline{\scriptsize\color{muted}code: \code{bool(res.get('converged'))}} \\
\code{state\_\allowbreak{}solve\_\allowbreak{}accepted} & FP64 state solve accepted & $A_\mathrm{res}$ & The final state solve was accepted on the true relative residual, recomputed in FP64, at the configured tolerance $\tau$. \newline{\scriptsize\color{muted}code: \code{rel is not None and rel <= self.cfg.tol}} \\
\code{no\_\allowbreak{}vram\_\allowbreak{}spill} & no VRAM spill & $A_\mathrm{no\,spill}$ & No growth of shared GPU memory beyond the spill limit was observed: the solve stayed in dedicated GPU memory. \newline{\scriptsize\color{muted}code: \code{not self.spill.spilled}} \\
\code{spill\_\allowbreak{}monitored} & spill monitored & $A_\mathrm{monitor}$ & Spill sampling ran during the iterations. Only the release rung enables it, so lower rungs fail this check and are never certified. \newline{\scriptsize\color{muted}code: \code{self.cfg.spill\_\allowbreak{}check\_\allowbreak{}every > 0}} \\
\code{volume\_\allowbreak{}ok} & volume within 0.1\% & $A_\mathrm{vol}$ & The largest recorded volume error is at most $10^{-3}$ (0.1\%). \newline{\scriptsize\color{muted}code: \code{res.get('max\_\allowbreak{}vol\_\allowbreak{}error') is not None and res['max\_\allowbreak{}vol\_\allowbreak{}error'] <= 0.001}} \\
\code{passive\_\allowbreak{}exact} & keep-in/out exact & $A_\mathrm{passive}$ & Every keep-in element has physical density exactly 1 and every keep-out element exactly 0. \newline{\scriptsize\color{muted}code: \code{bool(pc.get('keep\_\allowbreak{}in\_\allowbreak{}ok',\allowbreak{} True) and pc.get('keep\_\allowbreak{}out\_\allowbreak{}ok',\allowbreak{} True))}} \\
\end{fixtable}

\end{table}

\begin{figure}[t]
\wideshot{09_release_surface.png}
\caption{A certified release of the bracket: the marching-cubes surface of the 8.2-million-element design (on the reduced
ladder), the result card on the right and every rung's time on the ladder.}
\label{fig:releasesurface}
\end{figure}

\paragraph{Why lower rungs are never certified.} Spill monitoring samples GPU memory every \SpillEvery\ iterations,
and only release-rung solves enable it. A sketch, preview or refine result therefore fails \emph{spill monitored} and is
shown as \gui{not certified}, even when every other check passes (Figure~\ref{fig:b06refine}). It is feedback, not a
release. Only a certified result on the release rung is a released design, and the card says so whenever the latest
result is on a lower rung.

\paragraph{What the spill check means.} On a desktop GPU, the driver can move memory into shared system memory when
the card fills up. A solve then still runs, but its timings no longer mean anything. The spill guard compares shared
GPU memory against its value at the start of the run and stops the solve when it grows past the limit. A passing
check means no spill was observed under that sampling policy.

\begin{notebox}[What the certificate does not say]
The certificate belongs to the density problem that was solved, on the reported grid. It is not a proof of a global
optimum. Extracting a surface does not transfer the certificate to a new mechanical analysis of that surface, and the
certificate says nothing about stress, buckling, fatigue or manufacturing. In the paper's gate ablation, a problem with
its load inside a keep-out region converged and passed all six numerical checks; only rule \rref{R008} stops it
\claim{C4.4}.
\end{notebox}

\section{Surface, voxels and stale results}\label{sec:stale}
\index{stale result}
A certified release is shown as its $\rho = 0.5$ isosurface, extracted with marching cubes (VTK Flying Edges). For
large grids the surface on screen is pooled for display, as the note under the buttons says; the exported STL is not.
\gui{Voxels} shows the raw density field (Figure~\ref{fig:voxels}).

Each result is tied to the hash of the specification it solved. Change the specification after a solve and the card
marks its result \gui{stale} (Figure~\ref{fig:stale}): the design on screen no longer matches the specification. Undo the
change and the badge clears.

\begin{figure}[t]
\centering
\begin{minipage}[b]{0.6\linewidth}\shot{12_voxels.png}\end{minipage}\hfill
\begin{minipage}[b]{0.33\linewidth}\shot{14_stale.png}\end{minipage}
\caption{Left: the same design shown as density voxels. Right: a stale result after the volume fraction was edited.}
\label{fig:voxels}\label{fig:stale}
\end{figure}

\section{Export}\label{sec:export}
\index{export}
\gui{Export STL + report} (or the tools \tool{export_stl} and \tool{export_report}) writes three files to the session's
\code{exports/} folder:
\begin{itemize}
  \item a closed binary STL of the full-resolution $\rho = 0.5$ surface, positions in the spec's length unit; the
  export also checks, for surfaces of up to 30 million triangles, that every edge is shared by exactly two triangles and records the enclosed volume;
  \item a Markdown report and the same report as JSON, with the specification and its hash, the rung and solver,
  iterations and compute time, the compliance, every certificate check, the STL record with its SHA-256, and the hash of
  the session log (Appendix~\ref{app:certificate}).
\end{itemize}
A design that is not a certified result is exported with \code{\_draft} in its file name.
